% !TEX root = main.tex
\section*{Abstract — 국문 요약}
생산·이전 이력이 누적된 뒤 원료의 인정 근거에 문제가 발견되면, 과거 입력의 확인과 함께 현재 남아 있는 후속 회계 객체의 사용을 제한하는 대응을 고려할 수 있다. 입력 연결을 숨기는 공급망 원장에서는 공개 검증 결과만으로 그 대상과 집행 식별자를 얻기 어렵다.
\reviewcomment{01 — 문제 제기}{
    ``후속 회계 객체'' 같은 내부 용어보다, Privacy를 유지하면서 제품의 출처와 후속 생산 이력을 추적해야 한다는 문제를 먼저 설명하면 좋겠습니다.
    \par\smallskip\textbf{수정 제안}\par
    공급망의 조달·생산 연결을 숨기는 원장에서는 공개된 기록만으로 제품의 생산 이력이나 특정 원료가 사용된 후속 제품을 추적하기 어렵다.
    따라서 민감한 공급망 정보를 일반에 공개하지 않으면서, 승인된 감사자가 제품의 출처와 후속 생산 이력을 조사할 수 있는 구조가 필요하다.
}

본 논문은 비공개 회계 전이, 승인된 이력 복원과 소비 차단을 연결하는 zkDPP 구성을 제시한다.
실제 부모 참조와 출력의 미래 소비 식별자를 전이 증명에 결합하고, 신뢰된 감사자가 복원한 전체 미소비 대상을 상태 변경에 연결한다.
\reviewcomment{02 — 제안하는 방법}{
    ``증명에 결합''과 ``상태 변경에 연결'' 대신, 추적 정보의 정확성을 검증하고 미소비 대상을 동결한다는 동작을 직접 설명하면 좋겠습니다.
    \par\smallskip\textbf{수정 제안}\par
    본 논문은 공급망 거래를 비공개로 검증하면서 승인된 감사자의 양방향 추적과 후속 사용 차단을 지원하는 zkDPP를 제안한다.
    각 거래의 영지식증명은 암호화된 추적 정보가 실제 입력·출력과 일치하는지 검증한다.
    감사자는 이를 복호화해 과거 생산 이력과 후속 기록을 추적하고, 사용 중단이 승인되면 기준 시점에 아직 소비되지 않은 후속 대상을 동결한다.
}

보안을 정확성·건전성·프라이버시로 정의하고 명시적 계약 아래 증명한다.
정확성은 정상 조건 아래 전이·추적·동결의 기대 결과를, 건전성은 잘못된 전이·추적 성공·차단 완료의 배제와 차단 유지 중 소비 거부를 다룬다.
\reviewcomment{03 — 보안 보장}{
    ``명시적 계약''의 뜻을 풀고, 정상적인 추적 성공이 아니라 잘못된 성공·완료 판정을 배제한다는 점을 명확히 해야 합니다.
    \par\smallskip\textbf{수정 제안}\par
    명시한 암호학적 가정과 원장·권한 모델 아래 정확성, 건전성, 프라이버시를 분석한다.
    정상 실행에서 거래·추적·동결이 기대한 결과를 내며, 잘못된 거래의 수락과 부정확한 추적·차단 완료 판정을 배제함을 증명한다.
    전체 대상의 차단 상태가 유지되는 동안에는 해당 대상의 추가 소비가 거부된다.
}

프라이버시는 두 키 암호화의 기밀성을 바탕으로, 무작위 오라클 모형에서 비밀·개방값 공개가 없는 수동 외부 관찰자에 한정한다.
\reviewcomment{04 — Privacy의 대상과 가정}{
    ``두 키 암호화'' 같은 구성 용어보다 누구에게 어떤 정보를 숨기는지 먼저 설명하되, 수동 외부 관찰자에 한정된 보장임은 유지해야 합니다.
    \par\smallskip\textbf{수정 제안}\par
    프라이버시는 참여자의 비밀이나 객체의 비공개 내용을 열 수 있는 값을 얻지 못한 수동 외부 관찰자를 대상으로 한다.
    기록 암호와 감사키 보호 암호의 기밀성을 가정하고, 무작위 오라클 모형에서 공개 기록으로부터 비공개 정보가 보호됨을 증명한다.
}

개별 동결 성공과 전체 완료가 다른 경합·시점 반례를 제시한다.
\reviewcomment{05 — 완료 판정}{
    ``경합·시점 반례'' 대신 동결 전에 대상이 소비될 수 있어 개별 동결 성공만으로 전체 차단을 보장하지 못한다는 상황을 설명하면 좋겠습니다.
    \par\smallskip\textbf{수정 제안}\par
    동결 전에 대상이 소비되거나 대상별 상태를 서로 다른 시점에 확인하는 사례를 통해, 개별 동결의 성공만으로 전체 대상의 차단 완료를 판단할 수 없음을 보인다.
}

대응은 원장 객체의 추가 소비에 한정하며 기발급 제품 주장 취소나 물리적 회수는 포함하지 않는다.
기호 사례 검사는 별도의 모델 근거이며, 구체 암호 방식의 가정 충족과 구현 적합성·성능은 검증하지 않았다.
\reviewcomment{06 — 한계와 평가 범위}{
    ``기호 사례 검사''가 무엇을 확인한 것인지 쉽게 설명하고, 구현·성능을 검증하지 않았다는 문장이 실제 논문의 평가 범위와 맞는지 확인해야 합니다.
    \par\smallskip\textbf{수정 제안}\par
    동결은 원장 객체의 추가 소비를 막는 조치이며, 이미 발급된 DPP Claim의 취소나 실물 제품의 회수를 수행하지 않는다.
    기호 모델에서 거래·추적·차단 완료 판정의 사례를 검사했지만, 구체 암호 방식의 가정 충족이나 실제 구현의 적합성·성능을 검증한 것은 아니다.
}


\section*{Abstract — English}
When the evidence supporting a material input is disputed after production and transfer, a response may need to restrict the remaining descendant accounting objects.
In a ledger that hides input links, public verification alone does not identify these objects and their enforcement identifiers.
This paper presents zkDPP, a construction connecting private supply-chain accounting, authorized provenance recovery, and prevention of further consumption.
Event proofs bind encrypted parent references and future spend identifiers to the actual inputs and outputs.
A trusted auditor uses the recovered identifiers to trace the complete unspent frontier and change its status.
Under explicit construction contracts, we prove correctness, soundness, and privacy.
Correctness covers expected transition, tracing, and freezing results under their normal execution conditions.
Soundness excludes invalid transitions, incorrect successful traces, false completion, and further consumption while the full frontier remains blocked.
Privacy uses a two-key confidentiality lemma and a random-oracle proof for a passive external observer without participant-secret or opening disclosure.
Counterexamples involving consumption races and observations at different times distinguish individual freeze success from completion of the entire response.
The construction restricts subsequent ledger consumption; it does not revoke issued product claims or enforce physical recall.
Symbolic examples provide separate model evidence; they do not establish conformance of a concrete cryptographic instantiation or implementation performance.

\textbf{Keywords:} Digital product passport, Zero-knowledge proof, Provenance tracing, Supply-chain accounting, Authorized freezing
